\item En física es importante usar aproximaciones matemáticas. (a) Demuestre que, para ángulos pequeños ($< 20^\circ$),
$$ \tan \alpha \approx \operatorname{sen} \alpha \approx \alpha = \frac{\pi \alpha'}{180^\circ} $$
donde $\alpha$ está en radianes y $\alpha'$ en grados. (b) Use una calculadora para encontrar el ángulo más grande para el que tan $\alpha$ se pueda aproximar a $\alpha$ con un error menor de $10.0$ por ciento.
